\documentclass[11pt,a4paper]{ctexart}
\usepackage{amsmath, amsthm, amssymb, graphicx,mathrsfs}
\usepackage{geometry}
\usepackage[colorlinks=true,hidelinks]{hyperref} % 目录添加超链接
\geometry{left=1.24cm, right=1.24cm, top=1.91cm, bottom=1.91cm}
\title{最优控制理论}
\author{W·D·Gaster}

\renewcommand{\d}{{\rm d}}
\newcommand{\e}{{\rm e}}
\renewcommand{\j}{{\rm j}}
\newcommand{\bm}{\symbfit}
\usepackage{booktabs}
\usepackage{unicode-math}   % 添加\symbfit以实现加粗斜体命令
\usepackage{pifont} %使用带圈数字序号
\usepackage{enumitem}%自定义列表序号

\graphicspath{{figures/}} % 指定图片路径
\usepackage{float}

\begin{document}

\maketitle
\thispagestyle{empty}
\newpage

\tableofcontents
\thispagestyle{empty}
\newpage

\setcounter{page}{1}
受控系统的数学模型：$\bm{\dot x}(t)=\bm{f}(\bm{x}(t),\bm{u}(t),t),\bm{x}(t_0)=\bm{x}_0,t\geqslant t_0$

边界条件：起始状态$\bm{x}(t_0)=\bm{x}_0$是已知的，而最终到达的末态时刻和状态都可能固定或自由，可以用末态约束条件$\bm{g}[\bm{x}(t_f),t_f]=0$或$\bm{h}[\bm{x}(t_f),t_f]\leqslant0$来表示。

目标集：满足末态条件的状态集合$\bm{M}=\{\bm{x}(t_f)|\bm{x}(t_f)\in\mathbb{R}^n,\bm{g}[\bm{x}(t_f),t_f]=0,\bm{h}[\bm{x}(t_f),t_f]\leqslant0\}$

容许控制：控制向量$\bm{u}$的各分量具有不同物理属性，实际的物理量受到客观条件的限制，由约束条件规定的点集称为控制域$\bm{U}$，
在控制域内取值的每一个分段连续的控制向量$\bm{u}(t)$称为容许控制，记为$\bm{u}(t)\in\bm{U}$。

性能指标：取决于设计要求，性能指标函数（泛函）$J=\phi(\bm{x}(t_f),t_f)+\int_{t_0}^{t_f}L(\bm{x}(t),\bm{u}(t),t)\d t$。其中，
\begin{itemize}
    \item $\phi(\bm{x}(t_f),t_f)$是末值项，反映系统的末态接近目标集程度的度量，称为末值型性能指标；
    \item $\int_{t_0}^{t_f}L(\bm{x}(t),\bm{u}(t),t)\d t$是过程项，表示系统状态和控制的要求，反映能量的消耗，称为积分型性能指标。
\end{itemize}

设被控对象的状态方程及初始条件为：$\bm{\dot x}(t)=\bm{f}(\bm{x}(t),\bm{u}(t),t),\bm{x}(t_0)=\bm{x}_0$，目标集为$\bm{M}$，求取一个容许控制$\bm{u}(t)\in\bm{U},t\in[t_0,t_f]$，
使受控系统由给定初始条件出发，在末端时刻$t_f>t_0$，将系统转移到目标集$\bm{M}$，并使性能指标$J$最小。
\section{变分法}
    \subsection{泛函与变分的基本概念}
        \subsubsection{泛函}
            泛函：设$\{y(x)\}$为给定的某类函数，对于这类函数中的每一个函数，都有某个数$J$对应，则称$J$为这类函数的泛函，记为$J=J[y(x)]$，函数类$\{y(x)\}$称作泛函定义域，$y(x)$称作函数的自变量或者宗量。
            性能指标泛函可以看成线性赋范空间中的某个子集到实数集的映射，即$y=J[\bm{x}],\bm{x}\in\mathbb{R}^n,y\in\mathbb{R}$，称$J[\bm{x}]$为$\mathbb{R}^n$到$\mathbb{R}$的泛函算子。

            连续泛函：$\forall\varepsilon>0,\exists\delta$，当$\|\bm{x}(t)-\bm{x}_0(t)\|<\delta,\|\bm{\dot x}(t)-\bm{\dot x}_0(t)\|<\delta,\cdots,\|\bm{x}^{(k)}(t)-\bm{x}_0^{(k)}(t)\|<\delta$时，
            有$|J[\bm{x}(t)]-J[\bm{x}_0(t)]|<\varepsilon$，则称$J[\bm{x}(t)]$是在$\bm{x}_0(t)$处具有$k$阶接近度的连续泛函。

            线性泛函：设$J[\bm{x}(t)]$为连续泛函，若对于任意常数$\alpha,\beta$和定义域内任何变量$\bm{x}_1(t),\bm{x}_2(t)$，
            都有$J[\alpha\bm{x}_1(t)+\beta\bm{x}_2(t)]=\alpha J[\bm{x}_1(t)]+\beta J[\bm{x}_2(t)]$，则称$J[\bm{x}]$为线性泛函。
        \subsubsection{变分与运算规则}
            宗量变分：$\delta\bm{x}:=\bm{x}(t)-\bm{x}_0(t),\forall\bm{x}(t),\bm{x}_0(t)\in\mathbb{R}^n$

            泛函变分：若线性赋范空间上的连续泛函$J[\bm{x}]$的增量可表示为$\Delta J[\bm{x}]=J[\bm{x}+\delta\bm{x}]-J[\bm{x}]=L[\bm{x},\delta\bm{x}]+r[\bm{x},\delta\bm{x}]$，
            其中$L[\bm{x},\delta\bm{x}]$为关于$\delta\bm{x}$的线性连续泛函，$r[\bm{x},\delta\bm{x}]$是关于$\delta\bm{x}$的高阶无穷小，则$\delta J:=L[\bm{x},\delta\bm{x}]$称为泛函的变分。

            设$J[\bm{x}]$是连续赋范空间上的连续泛函，若在$\bm{x}=\bm{x}_0$处$J[\bm{x}]$可微，
            则$\delta J[\bm{x},\delta\bm{x}]=\dfrac{\partial}{\partial\varepsilon}J[\bm{x},\varepsilon\delta\bm{x}]\Bigg|_{\varepsilon\rightarrow0},0\leqslant\varepsilon\leqslant1$

            变分规则：
            \begin{itemize}
                \item $\delta(J_1+J_2)=\delta J_1+\delta J_2$
                \item $\delta(J_1J_2)=J_1\delta J_2+J_2\delta J_1$
                \item $\delta\int_a^bJ[\bm{x},\bm{\dot x},t]\d t=\int_a^b\delta J[\bm{x},\bm{\dot x},t]\d t$
                \item $\delta\dfrac{\d\bm{x}}{\d t}=\dfrac{\d}{\d t}\delta\bm{x}$
            \end{itemize}
        \subsubsection{泛函极值及判定}
            泛函极值：若对任意一个与$\bm{x}_0(t)$接近的变量$\bm{x}(t)$，都有$J[\bm{x}_0(t)]-J[\bm{x}(t)]\geqslant0$，则称泛函$J[\bm{x}]$在$\bm{x}_0$达到一个相对极大值。
            如果上述定义对定义域中所有$\bm{x}$都成立，称达到定义域上的绝对极大值。如果不等号反向，则称达到极小值。

            其中，$\bm{x}(t)$与$\bm{x}_0(t)$接近有两种含义：其一是$\|\bm{x}(t)-\bm{x}_0(t)\|$很小，具有零阶接近度，此时定义的极值称为强极值；
            其二$\|\bm{x}(t)-\bm{x}_0(t)\|$和$\|\bm{\dot x}(t)-\bm{\dot x}_0(t)\|$都很小，具有一阶接近度，此时定义的极值称为弱极值。

            极值原理：如果具有变分的泛函$J[\bm{x}]$在$\bm{x}=\bm{x}_0$上达到极值，则$J[\bm{x}]$沿着$\bm{x}_0$的变分$\delta J=0$。

            引理：设$\bm{M}(t)$为$[t_0,t_f]$上的r维连续向量函数，如果对于任意在$[t_0,t_f]$上连续且满足$\bm{y}(t_0)=\bm{y}(t_f)=0$的r维连续向量函数$\bm{y}(t)$，
            有$\int_{t_0}^{t_f}\bm{y}^T(t)\bm{M}(t)\d t=0$，则$\bm{M}(t)$在$[t_0,t_f]$上恒为零。
    \subsection{欧拉-拉格朗日方程}
        泛函$J[\bm{x}(t)]$在$\bm{x}^*(t)$取得极值的必要条件是：\newline        
        欧拉-拉格朗日方程
        \[
        \dfrac{\partial\phi}{\partial\bm{x}}-\dfrac{\d}{\d t}\dfrac{\partial\phi}{\partial\bm{\dot x}}=0
        \]
        和横截条件
        \[
        \bigg(\dfrac{\partial\phi}{\partial\bm{\dot x}}\bigg)^T\Bigg|_{t_f}\delta\bm{x}(t_f)-\bigg(\dfrac{\partial\phi}{\partial\bm{\dot x}}\bigg)^T\Bigg|_{t_0}\delta\bm{x}(t_0)=0
        \]
        对于最优解$\bm{x}=\bm{x}^*(t)$成立。

        (1)欧拉方程是一个时变二阶非线性微分方程，求解所需的两点边界条件由横截条件提供，在初值和终端值固定的情况下，$\delta\bm{x}(t_0)=0,\delta\bm{x}(t_f)=0$，则退化为边界条件$\bm{x}(t_0)=\bm{x}_0,\bm{x}(t_f)=\bm{x}_f$。

        (2)欧拉方程是泛函极值的必要条件，对于实际问题需要判断出极值的存在性再利用欧拉方程求出极值。

        泛函在约束$\bm{F}[\bm{x},\bm{\dot x},t]=0,t\in[t_0,t_f]$下的条件极值问题等价于
        \[
        J[\bm{x},\bm{\lambda}]=\int_{t_0}^{t_f}[\phi(\bm{x},\bm{\dot x},t)+\bm{\lambda}^T\bm{F}[\bm{x},\bm{\dot x},t]]\d t
        \]
        的无条件极值问题。即在约束$\bm{F}[\bm{x},\bm{\dot x},t]=0,t\in[t_0,t_f]$下$J[\bm{x}]$的最优轨线$\bm{x}^*(t)$上，泛函$J[\bm{x},\lambda]$的欧拉方程成立。
    \subsection{横截条件}
        求解欧拉-拉格朗日方程需要横截条件提供两点边界值$(t_0,\bm{x}_0)$和$(t_f,\bm{x}_f)$。

        末端固定时：末端时刻$t_f$固定，则横截条件一般表达式是
        \[
        \bigg(\dfrac{\partial\phi}{\partial\bm{\dot x}}\bigg)^T\Bigg|_{t_f}\delta\bm{x}(t_f)-\bigg(\dfrac{\partial\phi}{\partial\bm{\dot x}}\bigg)^T\Bigg|_{t_0}\delta\bm{x}(t_0)=0
        \]
        如果$\bm{x}(t_0)=\bm{x}_0$固定，则横截条件和边界条件为$\bigg(\dfrac{\partial\phi}{\partial\bm{\dot x}}\bigg)^T\Bigg|_{t_f}\delta\bm{x}(t_f)=0,\bm{x}(t_0)=\bm{x}_0$。

        如果初值和末端值固定，则横截条件和边界条件退化为$\bm{x}(t_f)=\bm{x}_f,\bm{x}(t_0)=\bm{x}_0$。

        如果末端时刻固定但状态自由，$\delta\bm{x}(t_f)\neq0$，条件变为$\dfrac{\partial\phi}{\partial\bm{\dot x}}\Bigg|_{t_f}=0,\bm{x}(t_0)=\bm{x}_0$。

        末端自由时：末端时刻自由，状态有约束$\bm{x}(t_f)=\bm{c}(t_f)$，则$\delta\bm{x}(t_f)=\delta{x}_f-\bm{\dot x}(t_f)\delta t_f,\delta\bm{x}_f=\bm{\dot c}(t_f)\delta t_f$
        $\bigg(\dfrac{\partial\phi}{\partial\bm{\dot x}}\bigg)^T\Bigg|_{t_f}\delta\bm{x}(t_f)+\phi(\bm{x}^*,\bm{\dot x}^*,t)\Bigg|_{t_f}\delta t_f=0,\bm{x}(t_0)=\bm{x}_0$
        
        若末端状态自由，则$\delta t_f=0$，演化为$\bigg(\dfrac{\partial\phi}{\partial\bm{\dot x}}\bigg)^T\Bigg|_{t_f}\delta\bm{x}(t_f)=0,\bm{x}(t_0)=\bm{x}_0$
        
        末端时刻和状态自由横截条件$\bm{x}_0=\bm{x}_0,\bigg(\dfrac{\partial\phi}{\partial\bm{\dot x}}\bigg)^T\Bigg|_{t_f}=0,[\phi(\bm{x}^*,\bm{\dot x}^*,t)-\bm{\dot x}^T(t)\dfrac{\partial\phi}{\partial\bm{\dot x}}]\Bigg|_{t_f}=0$，
        
        末端状态受约束：$\bm{x}_0=\bm{x}_0,\bm{x}(t_f)=\bm{c}(t_f),\bigg(\dfrac{\partial\phi}{\partial\bm{\dot x}}\bigg)^T\Bigg|_{t_f}
        =0,[\phi(\bm{x}^*,\bm{\dot x}^*,t)-(\bm{\dot c-\dot x})^T\dfrac{\partial\phi}{\partial\bm{\dot x}}]\Bigg|_{t_f}=0$
        
        初始端状态可变$\bigg(\dfrac{\partial\phi}{\partial\bm{\dot x}}\bigg)\Bigg|_{t_f}=0,\bigg(\dfrac{\partial\phi}{\partial\bm{\dot x}}\bigg)\Bigg|_{t_0}=0$
    \subsection{变分法在最优控制中的应用}
        n维系统的状态方程为$\bm{\dot x}(t)=\bm{f}(\bm{x}(t),\bm{u}(t),t)$，系统初始状态为$\bm{x}(t_0)=\bm{x}_0$，寻找一个无约束的连续的控制向量$\bm{u}(t)\in\mathbb{R}^r$，
        使得系统轨线在某一时刻$t_f>t_0$转移到目标集$\varPsi(\bm{x}(t_f),t_f)=0$里，其中$\varPsi\in\mathbb{R}^m,m\leqslant n$，
        使性能指标$J=\Phi(\bm{x}(t_f),t_f)+\int_{t_0}^{t_f}L(\bm{x},\bm{u},t)\d t$为最小。

        构造哈密顿函数$H(\bm{x},\bm{u},\bm{\lambda})=L(\bm{x},\bm{u},t)+\bm{\lambda}^T(t)\bm{f}(\bm{x},\bm{u},t)$，其中$\bm{\lambda}\in\mathbb{R}^n$为拉格朗日乘子向量。
        \subsubsection{末端时刻固定时最优解}
            构造广义泛函
            \begin{align*}
            J_a&=\Phi(\bm{x}(t_f))+\bm{\gamma}^T\varPsi(\bm{x}(t_f))+\int_{t_0}^{t_f}(L(\bm{x},\bm{u},t)+\bm{\lambda}^T[\bm{f}(\bm{x},\bm{u},t)-\bm{\dot x}])\d t\\
            &=\Phi(\bm{x}(t_f))+\bm{\gamma}^T\varPsi(\bm{x}(t_f))+\int_{t_0}^{t_f}[H(\bm{x},\bm{u},\bm{\lambda},t)-\bm{\lambda}^T(t)\bm{\dot x}(t)]\d t\\
            &=\Phi(\bm{x}(t_f))+\bm{\gamma}^T\varPsi(\bm{x}(t_f))-\bm{\lambda}^T(t_f)\bm{x}(t_f)+\bm{\lambda}^T(t_0)\bm{x}(t_0)+\int_{t_0}^{t_f}[H(\bm{x},\bm{u},\bm{\lambda},t)+\bm{\dot\lambda}^T(t)\bm{x}(t)]\d t\\
            \end{align*}

            对广义泛函取变分（$\bm{\lambda}(t)$和$\bm{\gamma}$不变分），得
            \[
            \delta J_a=\delta\bm{x}^T(t_f)\bigg[\dfrac{\partial\Phi}{\partial\bm{x}(t_f)}+\dfrac{\partial\varPsi^T}{\partial\bm{x}(t_f)}\bm{\gamma}-\bm{\lambda}(t_f)\bigg]+
            \int_{t_0}^{t_f}\bigg[\bigg(\dfrac{\partial H}{\partial\bm{x}}+\bm{\dot\lambda}\bigg)^T\delta\bm{x}+\bigg(\dfrac{\partial H}{\partial\bm{u}}\bigg)^T\delta\bm{u}\bigg]\d t
            \]
            令$\delta J_a=0$，由于$\delta\bm{x},\delta\bm{u},\delta\bm{x}(t_f)$的任意性，得必要条件为\newline
            欧拉方程：
            \[
            \bm{\dot\lambda}(t)=-\dfrac{\partial H}{\partial\bm{x}}
            \dfrac{\partial H}{\partial\bm{u}}=0
            \]
            横截条件：
            \[
            \bm{\lambda}(t_f)=\dfrac{\partial\Phi}{\partial\bm{x}(t_f)}+\dfrac{\partial\varPsi^T}{\partial\bm{x}(t_f)}\bm{\gamma}
            \]
            加上引入哈密顿函数的条件$\bm{\dot x}(t)=\dfrac{\partial H}{\partial\bm{\lambda}}=\bm{f}(\bm{x},\bm{u},t)$

            对于如下最优控制问题：$\min\limits_{\bm{u}(t)}J=\Phi(\bm{x}(t_f),t_f)+\int_{t_0}^{t_f}L(\bm{x},\bm{u},t)\d t$

            使得$\bm{\dot x}(t)=\bm{f}(\bm{x},\bm{u},t),\bm{x}(t_0)=\bm{x}_0,\varPsi(\bm{x}(t_f))=0$
            最优解的必要条件为：
            
            (1)正则方程：$\bm{\dot\lambda}(t)=-\dfrac{\partial H}{\partial\bm{x}},\bm{\dot x}(t)=\dfrac{\partial H}{\partial\bm{\lambda}}=\bm{f}(\bm{x},\bm{u},t)$
            
            (2)边界条件与横截条件：$\bm{x}(t_0)=\bm{x}_0,\varPsi(\bm{x}(t_f))=0,\bm{\lambda}(t_f)=\dfrac{\partial\Phi}{\partial\bm{x}(t_f)}+\dfrac{\partial\varPsi^T}{\partial\bm{x}(t_f)}\bm{\gamma}$
            
            (3)极值条件（控制方程）：$\dfrac{\partial H}{\partial\bm{u}}=0$

            (1)当末端状态自由时，由于不存在目标集，因此去除结论中带有$\varPsi(\bm{x}(t_f))$的部分即可得到时刻固定状态自由的泛函极值必要条件
            
            (2)当末端状态时刻均固定时，广义泛函的一次变分变为：$\delta J_a=\int_{t_0}^{t_f}\bigg[\bigg(\dfrac{\partial H}{\partial\bm{x}}+\bm{\dot\lambda}\bigg)^T\delta\bm{x}
            +\bigg(\dfrac{\partial H}{\partial\bm{u}}\bigg)^T\delta\bm{u}\bigg]\d t$，
            其中$\delta\bm{x}$任意而$\delta\bm{u}$必须满足某些限制条件，即得到欧拉方程，此时横截条件不成立，$\bm{x}(t_0)=\bm{x}_0$和$\bm{x}(t_f)=\bm{x}_f$是求解正则方程的两端边界条件。
        \subsubsection{末端时刻自由时的最优解}
            在上述条件基础上，增加第四个条件：在最优轨线的末端哈密顿函数的变化律$H(t_f)=-\dfrac{\partial\Phi}{\partial t_f}-\bm{\gamma}^T\dfrac{\partial\varPsi}{\partial t_f}$

            对于末端状态和时刻都自由的情形，正则方程不变，横截条件和边界条件变为：$\bm{x}(t_0)=\bm{x}_0\qquad\bm{\lambda}(t_f)=\dfrac{\partial\Phi}{\partial\bm{x}(t_f)}$，
            哈密顿函数变化律为$H(t_f)=-\dfrac{\partial\Phi}{\partial t_f}$
            
            对于末端状态固定的情形，正则方程不变，横截条件退化成两点边界条件，哈密顿函数变化律为$H(t_f)=-\dfrac{\partial\Phi}{\partial t_f}$
\newpage
\section{极小值原理}
    古典变分法要求控制向量不受约束，要求$\partial H/\partial\bm{u}$存在且连续，当容许控制集合是闭集时边界上$\delta\bm{u}$不能任意无法满足$\partial H/\partial\bm{u}=0$。
    为解决控制有约束的变分问题，Pontryagin提出极小值原理，Bellman提出动态规划，能够应用于受边界限制的情形，并且不必要要求哈密顿函数对控制向量连续可微。
    \subsection{连续系统末端自由的极小值原理}
        对于如下定常系统、末值型性能指标、末端状态自由、控制受约束的最优控制问题：
        \[
        \min\limits_{\bm{u}(t)\in\bm{\Omega}}J(\bm{u})=\Phi(\bm{x}(t_f))
        s.t.\bm{\dot x}=\bm{f}(\bm{x},\bm{u}),\bm{x}(t_0)=\bm{x}_0,\varPsi(\bm{x}(t_f))=0
        \]
        其中，$\bm{x}(t)\in\mathbb{R}^n,\bm{u}(t)\in\bm{\Omega}\subset\mathbb{R}^r$为任意分段连续函数，$\bm{\Omega}$为容许控制域，末端状态自由。并假设：

        (1)向量函数$\bm{f}(\bm{x},\bm{u})$和标量函数$\Phi(\bm{x})$都是连续可微函数；
        
        (2)在有界集上$\bm{f}(\bm{x},\bm{u})$对$\bm{x}$满足李普希兹条件$|\bm{f}(\bm{x}^1,\bm{u})-\bm{f}(\bm{x}^2,\bm{u})|\leqslant a|\bm{x}^1-\bm{x}^2|,a>0$
        
        则对于最优解$\bm{u}^*(t),t_f^*,\bm{x}^*(t)$，必存在非零的向量$\bm{\lambda}(t)\in\mathbb{R}^n$，使以下必要条件成立：
        (1)正则方程：$\bm{\dot\lambda}=-\dfrac{\partial H}{\partial\bm{x}},\bm{\dot x}(t)=\dfrac{\partial H}{\partial\bm{\lambda}}=\bm{f}(\bm{x},\bm{u})$
        其中哈密顿函数$H(\bm{x},\bm{u},\bm{\lambda})=\bm{\lambda}^T(t)\bm{f}(\bm{x},\bm{u})$
        
        (2)边界条件与横截条件：$\bm{x}(t_0)=\bm{x}_0,\bm{\lambda}(t_f)=\dfrac{\partial\Phi}{\partial\bm{x}(t_f)}$
        
        (3)极小值条件：$H(\bm{x}^*,\bm{u}^*,\bm{\lambda})=\min\limits_{\bm{u}\in\bm{\Omega}}H(\bm{x}^*,\bm{u},\bm{\lambda})$
        
        (4)沿最优轨线变化的哈密顿函数变化律（$t_f$自由时）：$H(\bm{x}^*(t_f^*),\bm{u}^*(t_f^*),\bm{\lambda}(t_f^*))=0$
        
        极小值原理将容许条件放宽了，对通常的控制约束均适用；
        最优控制使哈密顿函数取全局极小值。经典变分法的极值条件是极小值原理中极小值条件的一种特例；
        极小值原理不要求哈密顿函数对控制向量的可微性，放宽了应用条件。
    \subsection{几种问题极小值原理的转化}
        \subsubsection{非定常系统}
            (1)正则方程：$\bm{\dot\lambda}(t)=-\dfrac{\partial H}{\partial\bm{x}},\bm{\dot x}(t)=\dfrac{\partial H}{\partial\bm{\lambda}}=\bm{f}(\bm{x},\bm{u})$，
            其中哈密顿函数$H(\bm{x},\bm{u},\bm{\lambda},t)=\bm{\lambda}^T(t)\bm{f}(\bm{x},\bm{u},t)$

            (2)边界条件与横截条件：$\bm{x}(t_0)=\bm{x}_0,\bm{\lambda}(t_f)=\dfrac{\partial\Phi}{\partial\bm{x}(t_f)}$

            (3)极小值条件：$H(\bm{x}^*,\bm{u}^*,\bm{\lambda},t)=\min\limits_{\bm{u}\in\bm{\Omega}}H(\bm{x}^*,\bm{u},\bm{\lambda},t)$

            (4)沿最优轨线变化的哈密顿函数变化律（$t_f$自由时）：$H(\bm{x}^*(t_f^*),\bm{u}^*(t_f^*),\bm{\lambda}(t_f^*),t_f^*)=-\dfrac{\partial\Phi(\bm{x}^*(t_f^*),t_f^*)}{\partial t_f}$

            证明方法：设$x_{n+1}(t)=1$，则$\dot{x}_{n+1}(t)=1,x_{n+1}(t_0)=t_0,x_{n+1}(t_f)=t_f$
            构造增广向量
            
            $\bm{\bar x}(t)=\begin{bmatrix}\bm{x}(t)\\x_{n+1}(t)\end{bmatrix},\bm{\bar f}(\bm{\bar x},\bm{u})=
            \begin{bmatrix}\bm{f}(\bm{x},\bm{u},t)\\1\end{bmatrix},\bm{\bar x}_0(t_0)=\begin{bmatrix}\bm{x}_0\\t_0\end{bmatrix}=\bm{\bar x}_0$
            则将非定常问题转化为定常问题。
        \subsubsection{积分型性能指标}
            性能指标泛函$J(\bm{u})=\int_{t_0}^{t_f}L(\bm{x},\bm{u})\d t$，其他同非定常系统。

            (1)正则方程：$\bm{\dot\lambda}(t)=-\dfrac{\partial H}{\partial\bm{x}},\bm{\dot x}(t)=\dfrac{\partial H}{\partial\bm{\lambda}}=\bm{f}(\bm{x},\bm{u})$
            其中哈密顿函数$H(\bm{x},\bm{u},\bm{\lambda})=L(\bm{x},\bm{u})+\bm{\lambda}^T(t)\bm{f}(\bm{x},\bm{u})$
            
            (2)边界条件与横截条件：$\bm{x}(t_0)=\bm{x}_0,\bm{\lambda}(t_f)=0$
            
            (3)极小值条件：$H(\bm{x}^*,\bm{u}^*,\bm{\lambda})=\min\limits_{\bm{u}\in\bm{\Omega}}H(\bm{x}^*,\bm{u},\bm{\lambda})$
            
            (4)沿最优轨线变化的哈密顿函数变化律（$t_f$自由时）：$H(\bm{x}^*(t_f^*),\bm{u}^*(t_f^*),\bm{\lambda}(t_f^*))=0$

            证明方法：设$x_0(t)=\int_{t_0}^{t}L(\bm{x},\bm{u})\d t$，则$\dot{x}_0(t)=L(\bm{x},\bm{u}),x_0(t_0)=0,x_0(t_f)=\int_{t_0}^{t_f}L(\bm{x},\bm{u})\d t=J(\bm{u})$，
            构造增广向量$\bm{\bar x}(t)=\begin{bmatrix}x_0(t)\\\bm{x}(t)\end{bmatrix},\bm{\bar f}(\bm{\bar x},\bm{u})=\begin{bmatrix}L(\bm{x},\bm{u})\\\bm{f}(\bm{x},\bm{u})\end{bmatrix},
            \bm{\bar x}(t_0)=\begin{bmatrix}0\\\bm{x}_0\end{bmatrix}=\bm{\bar x}_0$
            则将非定常问题转化为定常问题。
        \subsubsection{末态约束}
            设末端状态约束为$\bm{g}[\bm{x}(t_f)]=0,\bm{h}[\bm{x}(t_f)]\leqslant0$
            若性能指标是末值型的，即$J(\bm{u})=\Phi[\bm{x}(t_f)]$，利用拉格朗日乘子法，则等价性能指标为
            $\bar J=\Phi[\bm{x}(t_f)]+\bm{\mu}^T\bm{g}(\bm{x}(t_f))+\bm{\nu}^T\bm{h}(\bm{x}(t_f))$，
            其中$\nu_i\geqslant0,\nu_ih_i(\bm{x}(t_f))=0$
            
            横截条件为$\bm{\lambda}(t_f)=\dfrac{\partial\Phi(\bm{x}(t_f))}{\partial\bm{x}(t_f)}+\dfrac{\partial\bm{g}^T(\bm{x}(t_f))}{\partial\bm{x}(t_f)}\bm{\mu}
            +\dfrac{\partial\bm{h}^T(\bm{x}(t_f))}{\partial\bm{x}(t_f)}\bm{\nu}$，
            其他条件不变。
        \subsubsection{有积分限制情形}
            对于如下定常系统、积分型性能指标、末端状态自由、控制和状态轨线受约束的最优控制问题：
            \[
            \min\limits_{\bm{u}(t)\in\bm{\Omega}}J(\bm{u})=\int_{t_0}^{t_f}L(\bm{x},\bm{u})\d t
            s.t.\bm{\dot x}=\bm{f}(\bm{x},\bm{u}),\bm{x}(t_0)=\bm{x}_0,\varPsi(\bm{x}(t_f))=0
            \]
            \[
            J_1=\int_{t_0}^{t_f}\bm{L}_1(\bm{x}(t),\bm{u}(t))\d t=0,J_2=\int_{t_0}^{t_f}\bm{L}_2(\bm{x}(t),\bm{u}(t))\d t\leqslant0
            \]

            则对于最优解的必要条件：

            (1)正则方程：$\bm{\dot\lambda}=-\dfrac{\partial H}{\partial\bm{x}},\bm{\dot x}(t)=\dfrac{\partial H}{\partial\bm{\lambda}}=\bm{f}(\bm{x},\bm{u})$
            其中哈密顿函数$H(\bm{x}(t),\bm{u}(t),\bm{\lambda}(t),\bm{\lambda}_1,\bm{\lambda}_2)=L(\bm{x},\bm{u})+\bm{\lambda}^T(t)\bm{f}(\bm{x},\bm{u})
            +\bm{\lambda}_1^T(t)\bm{L}_1(\bm{x},\bm{u})+\bm{\lambda}_2^T(t)\bm{L}_2(\bm{x},\bm{u})$，满足$\lambda_{2i}\geqslant0,\lambda_{2i}J_{2i}=0,i\leqslant l$
            
            (2)边界条件与横截条件：$\bm{x}(t_0)=\bm{x}_0\qquad\bm{\lambda}(t_f)=0$
            
            (3)极小值条件：$H(\bm{x}^*(t),\bm{u}^*(t),\bm{\lambda}(t),\bm{\lambda}_1,\bm{\lambda}_2)=\min\limits_{\bm{u}\in\bm{\Omega}}H(\bm{x}^*,\bm{u},\bm{\lambda},\bm{\lambda}_1,\bm{\lambda}_2)$
            
            (4)沿最优轨线变化的哈密顿函数变化律：
            
            $t_f$自由时，$H(\bm{x}^*(t_f^*),\bm{u}^*(t_f^*),\bm{\lambda}(t_f^*),\bm{\lambda}_1,\bm{\lambda}_2)=0$

            $t_f$固定时，$H(\bm{x}^*(t),\bm{u}^*(t),\bm{\lambda}(t),\bm{\lambda}_1,\bm{\lambda}_2)=H(\bm{x}^*(t_f^*),\bm{u}^*(t_f^*),\bm{\lambda}(t_f^*),\bm{\lambda}_1,\bm{\lambda}_2)={\rm const}$

            证明方法：扩充状态变量法，$\dot{z}_0=L(\bm{x},\bm{u}),z_0(t_0)=0,\dot{z}_1=L_1(\bm{x},\bm{u}),z_1(t_0)=0,\dot{z}_2=L_2(\bm{x},\bm{u}),z_2(t_0)=0$，
            从而，$J=z_0(t_f),J_1=z_1(t_f),J_2=z_2(t_f)$
            构造增广向量$\bm{\bar x}=\begin{bmatrix}z_0\\\bm{x}\\z_1\\z_2\end{bmatrix},\bm{\bar f}=\begin{bmatrix}L\\\bm{f}\\L_1\\L_2\end{bmatrix},\bm{\bar x}_0=\begin{bmatrix}0\\\bm{x}_0\\0\\0\end{bmatrix}$
            从而转化为具有末态约束、末值型性能指标的定常问题。
    \subsection{离散系统的极小值原理}
        $\bm{x}(k+1)=\bm{f}(\bm{x}(k),\bm{u}(k),k),\bm{x}(0)=\bm{x}_0$，性能指标$J=\Phi(\bm{x}(N))+\sum\limits_{k=0}^{N-1}L[\bm{x}(k),\bm{u}(k),k]$
        \subsubsection{末端约束时}
            末端状态受$\varPsi[\bm{x}(N)]=0$约束。则最优解满足：

            (1)差分方程：$\bm{x}(k+1)=\dfrac{\partial H(k)}{\partial\bm{\lambda}(k+1)},\bm{\lambda}(k)=\dfrac{\partial H(k)}{\partial\bm{x}(k)}$，
            其中离散哈密顿函数为$H(k)=H[\bm{x}(k),\bm{u}(k),k]=L[\bm{x}(k),\bm{u}(k),k]+\bm{\lambda}^T(k+1)\bm{f}[\bm{x}(k),\bm{u}(k),k]$
            
            (2)边界条件和横截条件：$\bm{x}(0)=\bm{x}_0,\bm{\varPsi}[\bm{x}(N)]=0,\bm{\lambda}(N)=\dfrac{\partial\Phi}{\partial\bm{x}(N)}+\dfrac{\partial\varPsi^T}{\partial\bm{x}(N)}\bm{\gamma}$
            
            (3)极小值条件：$H^*(k)=\min\limits_{\bm{u}(k)\in\bm{\Omega}}H(k)$，若$\bm{u}(k)$无约束，则为$\dfrac{\partial H(k)}{\partial\bm{u}(k)}=0$
        \subsubsection{末端自由时}
            末端状态自由。则最优解满足：
            
            (1)差分方程：$\bm{x}(k+1)=\dfrac{\partial H(k)}{\partial\bm{\lambda}(k+1)},\bm{\lambda}(k)=\dfrac{\partial H(k)}{\partial\bm{x}(k)}$，
            其中离散哈密顿函数为$H(k)=H[\bm{x}(k),\bm{u}(k),k]=L[\bm{x}(k),\bm{u}(k),k]+\bm{\lambda}^T(k+1)\bm{f}[\bm{x}(k),\bm{u}(k),k]$
            
            (2)边界条件和横截条件：$\bm{x}(0)=\bm{x}_0,\bm{\lambda}(N)=\dfrac{\partial\Phi}{\partial\bm{x}(N)}$
            
            (3)极小值条件：$H^*(k)=\min\limits_{\bm{u}(k)\in\bm{\Omega}}H(k)$，若$\bm{u}(k)$无约束，则为$\dfrac{\partial H(k)}{\partial\bm{u}(k)}=0$
    \subsection{能量最小控制问题}
        设线性定常系统状态方程为$\bm{\dot x}(t)=\bm{A}\bm{x}(t)+\bm{B}\bm{u}(t),\bm{x}(t_0)=\bm{x}_0$，控制约束$|u_i(t)|\leqslant M$，末端状态$\bm{x}(t_f)=\bm{x}_f$，确定最优控制使性能指标
        $J=\int_{t_0}^{t_f}\bm{u}^T(t)\bm{u}(t)\d t=\int_{t_0}^{t_f}\sum u_i^2(t)\d t$

        $H=\bm{u}^T(t)\bm{u}(t)+\bm{\lambda}^T(t)\bm{A}\bm{x}(t)+\bm{\lambda}^T(t)\bm{B}(t)\bm{u}(t)$
        定义开关向量函数$\bm{s}(t)=\bm{B}^T\bm{\lambda}(t)$

        由协态方程$\bm{\dot\lambda}(t)=-\dfrac{\partial H}{\partial\bm{x}}=-\bm{A}^T\bm{\lambda}(t)$解得$\bm{\lambda}(t)=\e^{-\bm{A}^T(t-t_0)}\bm{\lambda}(t_0)$，
        故开关函数$\bm{s}(t)=\bm{B}^T\e^{-\bm{A}^T(t-t_0)}\bm{\lambda}(t_0)$
        则哈密顿函数表示为$H=\sum[u_j^2(t)+u_j(t)s_j(t)]+\bm{x}^T(t)\bm{A}^T\bm{\lambda}(t)$

        控制向量没有约束则$\dfrac{\partial}{\partial\bm{u}_j(t)}[u_j^2(t)+u_j(t)s_j(t)]=0$，得$u_j^*(t)=-\dfrac{1}{2}s_j(t)$，由于$|u_j(t)|\leqslant M$，因此当$s_j(t)>2M$时，$u^*_j(t)=-M{\rm sign}[s_j(t)]$
    \subsection{时间最优控制问题}
        设线性定常系统状态方程为$\bm{\dot x}(t)=\bm{f}(\bm{x}(t),t)+\bm{B}(\bm{x}(t),t)\bm{u}(t)$，控制约束$|u_i(t)|\leqslant 1$，初态$\bm{x}(t_0)=\bm{x}_0$，末态$\bm{g}(\bm{x}(T),T)=0$，确定最优控制使性能指标
        $J=\int_{t_0}^{T}\d t=T-t_0$

        Bang-Bang控制原理：设$\bm{u}^*(t)$时时间最优控制，$\bm{x}^*(t)$和$\bm{\lambda}^*(t)$是相应的状态和协态。若问题是正常的，即对于所有$q_j(t)$只在$[t_0,T]$上有限个孤立点上为零，则对所有的$t\in[t_0,T]$，有
        \[
        u^*_j(t)=-{\rm sgn}[b_j^T(x^*(t),t)\lambda(t)]
        \]
        或者
        \[
        \bm{u}^*(t)=-{\rm sgn}[\bm{B}^T(x^*(t),t)\bm{\lambda}(t)]
        \]
        成立，即时间控制各分量都是时间的分段常值函数，且仅在开关$q(t)=0$的开关时间$t_{\beta_j}$上发生由一个恒值向另一个恒值的跳变。

        问题正常性的判定：r个矩阵$\bm{G}_i=[\bm{b}_i\quad\bm{Ab}_i\quad\bm{A^2b}_j\quad\cdots\quad\bm{A^{n-1}b}_i]$只要有一个是奇异的，问题就是奇异的，即正常问题要求每个$[\bm{A},\bm{b}_i]$是能控对，等价于受控系统对每个控制分量都是完全能控的。
\newpage
\section{线性二次型最优控制问题（LQR）}
    设线性定常系统受控对象状态方程为
    \[
    \begin{cases}
    \bm{\dot x}=\bm{A}(t)\bm{x}+\bm{B}(t)\bm{u}\\\bm{y}=\bm{C}(t)\bm{x},\bm{x}(t_0)=\bm{x}_0
    \end{cases}
    \]

    性能指标为$J=\dfrac12\bm{e}^T(t_f)\bm{Se}(t_f)+\dfrac12\int_{t_0}^{t_f}[\bm{e}^T(t)\bm{Q}(t)\bm{e}(t)+\bm{u}^T(t)\bm{R}(t)\bm{u}(t)]\d t$

    其中$\bm{e}(t)=\bm{x}(t)-\bm{x_r}(t)$为状态跟踪误差向量，$\bm{Q}(t)$为正半定状态加权矩阵，$\bm{R}(t)$为正定控制加权矩阵，$\bm{S}$为半正定终端加权矩阵。

    若取$\bm{S}={\rm diag}(s_1,s_2,\cdots,s_l)$，则末值项$\dfrac12\bm{e}^T(t_f)\bm{Se}(t_f)=\dfrac12\sum\limits_{i=1}^ls_i\bm{e}_i^2(t_f)$，
    即末态跟踪误差向量各分量的加权平方和，表示跟踪结束后对系统末态跟踪误差的要求。

    若取$\bm{Q}(t)={\rm diag}(q_1(t),q_2(t),\cdots,q_l(t))$，则积分项$\dfrac12\int_{t_0}^{t_f}\bm{e}^T(t)\bm{Q}(t)\bm{e}(t)\d t=\dfrac12\sum\limits_{i=1}^l\int_{t_0}^{t_f}q_i(t)e_i^2(t)\d t$，
    即跟踪过程中系统跟踪误差向量的加权平方和的积分，作为控制过程动态特性的一种度量。

    若取$\bm{R}(t)={\rm diag}(r_1(t),r_2(t),\cdots,r_l(t))$，则积分项$\dfrac12\int_{t_0}^{t_f}\bm{u}^T(t)\bm{R}(t)\bm{u}(t)\d t=\dfrac12\sum\limits_{i=1}^l\int_{t_0}^{t_f}r_i(t)u_i^2(t)\d t$，
    即控制过程中系统所消耗的控制能量。
    \subsection{状态调节器}
        令$\bm{C}(t)=\bm{I},\bm{x}_r(t)=0$，则$\bm{e}(t)=\bm{y}(t)=\bm{x}(t)$

        性能指标变为$J=\dfrac12\bm{x}^T(t_f)\bm{Sx}(t_f)+\dfrac12\int_{t_0}^{t_f}[\bm{x}^T(t)\bm{Q}(t)\bm{x}(t)+\bm{u}^T(t)\bm{R}(t)\bm{u}(t)]\d t$
        
        当系统受扰动影响偏离了原零平衡状态，要求控制向量使得系统状态恢复到原平衡点附近，并使性能指标极小。
        \subsubsection{有限时间状态调节器}
            在满足受控对象状态方程的约束条件下，在容许控制的范围内确定最优控制$\bm{u}^*(t)$。使得系统的状态在限定时间$[t_0.t_f]$内由给定的$\bm{x}_0$转移到$\bm{x}(t_f)$且使性能指标取极小值。

            最优状态调节器为：$\bm{u}(t)=\bm{K}(t)\bm{x}(t),\bm{K}(t)=-\bm{R}^{-1}(t)\bm{B}^T(t)\bm{P}(t)$，
            其中$\bm{P}(t)$满足矩阵微分Riccati方程
            \[
            -\bm{\dot P}(t)=\bm{P}(t)\bm{A}(t)+\bm{A}^T(t)\bm{P}(t)-\bm{P}(t)\bm{B}(t)\bm{R}^{-1}(t)\bm{B}^T(t)\bm{P}(t)+\bm{Q}(t)
            \]
            边界条件$\bm{P}(t_f)=\bm{S}$，且最优性能值为$J^*[\bm{x}(t)]=\dfrac12\bm{x}^T(t_0)\bm{P}(t_0)\bm{x}(t_0)$
        \subsubsection{无限长时间定常状态调节器}
            系统受扰动偏离原平衡状态后，能最优地恢复到平衡状态，不产生稳态误差。这类问题采用无限时间状态调节器。

            此时有$\bm{P}(t_f)=\bm{S}=0$，且末端项为零，令$t_0=0$，则性能指标$J=\dfrac12\int_{t_0}^{t_f}[\bm{x^T Qx}+\bm{u^T Ru}]\d t$，其中$\bm{R}$为常数对称正定矩阵，$\bm{Q}$为常数对称半正定矩阵。

            若$[\bm{A\quad B}]$是可控的，$\bm{R}>0$，$\bm{Q}=\bm{D^T D}\geqslant$，且$[\bm{A\quad D}]$可观，则Riccati方程满足$\bm{P}(t_f)=0$的解在无穷时存在极限且是唯一的常数矩阵$\lim\limits_{t_f\rightarrow\infty}\bm{P}(t_f)=\bm{P}$
            该极限为Riccati代数方程$\bm{PA}+\bm{A^T P}-\bm{PBR^{-1}B^T P}+\bm{Q}=0$的唯一正定解。

            此时存在唯一最优控制$\bm{u}^*(t)=-\bm{R^{-1}B^T PX}(t)$，最优性能指标$J^*[\bm{x}^*(t_0)]=\dfrac12\bm{x}^{*T}(t_0)\bm{Px}^*(t_0)$，对应的最优轨迹为$\bm{\dot x}^*(t)=[\bm{A-BR^{-1}B^T P}]\bm{x}^*(t)$。

            此时无限时间的二次型最优状态调节器控制系统的闭环系统是渐进稳定的，即闭环系统的系统矩阵$[\bm{A-BR^{-1}B^T P}]$具有负实部的特征值。
    \subsection{输出调节器}
        当系统受到外部扰动时，系统的实际输出能够在一定的性能指标下最优地恢复到原输出平衡状态或其邻近，称作输出调节器问题，是状态调节器的一种特殊情形。
        \subsubsection{线性时变系统}
        性能指标$J=\dfrac12\bm{y}^T(t_f)\bm{Sy}(t_f)+\dfrac12\int_{t_0}^{t_f}[\bm{y}^T\bm{Q}(t)\bm{y}+\bm{u}^T\bm{R}(t)\bm{u}]\d t$，
        
        最优输出调节器$\bm{u}^*(t)=-\bm{R^{-1}}(t)\bm{B}^T(t)\bm{P}(t)\bm{x}(t)$，其中$\bm{P}(t)$满足Riccati方程
        \[
        -\bm{\dot P}(t)+\bm{P}(t)\bm{A}(t)+\bm{A}^T(t)\bm{P}(t)-\bm{P}(t)\bm{B}(t)\bm{R}^{-1}(t)\bm{B}^T(t)\bm{P}(t)+\bm{C}^T(t)\bm{Q}(t)\bm{C}(t)=0
        \]

        只需将$\bm{y}=\bm{C}(t)\bm{x}$代入性能指标即可转化为一般状态调节器的性能指标进行计算。
        \subsubsection{线性定常系统}
        性能指标$J=\dfrac12\int_{t_0}^{t_f}[\bm{y^T Qy}+\bm{u^T Ru}]\d t$，
        
        最优输出调节器$\bm{u}^*(t)=-\bm{R^{-1}B^T Px}(t)$，其中$\bm{P}(t)$是Riccati代数方程
        \[
        \bm{PA}+\bm{A^T P}-\bm{PBR^{-1}B^T P}+\bm{C^T QC}=0
        \]
        的唯一正定解，在此作用律下系统是渐进稳定的。
    \subsection{最优跟踪问题}
        确定最优控制律使得系统输出$\bm{y}(t)$尽可能逼近某一要求的输出参考信号$\bm{z}(t)$，并使性能指标为最小。
        \subsubsection{线性时变系统}
            性能指标$J=\dfrac12(\bm{z}(t_f)-\bm{y}(t_f))^T\bm{S}(\bm{z}(t_f)-\bm{y}(t_f))+\dfrac12\int_{t_0}^{t_f}[(\bm{z-y})^T\bm{Q}(t)(\bm{z-y})+\bm{u}^T\bm{R}(t)\bm{u}]\d t$

            最优控制$\bm{u}^*(t)=-\bm{R^{-1}}(t)\bm{B}^T(t)\bm{P}(t)\bm{x}(t)+\bm{R}^{-1}(t)\bm{B}^T(t)\bm{g}(t)$，其中$\bm{P}(t)$和$\bm{g}(t)$满足
            \[
            \begin{cases}
            \bm{\dot P}(t)=\bm{P}(t)\bm{A}+\bm{A}^T\bm{P}(t)-\bm{P}(t)\bm{BR}^{-1}\bm{B}^T\bm{P}(t)+\bm{C}^T\bm{Q}(t)\bm{C}\\
            \bm{\dot g}(t)=[\bm{P}(t)\bm{BR^{-1}B^T-A^T}]\bm{g}(t)+\bm{C^T Qz}(t)
            \end{cases}
            \]

            边界条件$\bm{P}(t_f)=\bm{C}^T(t_f)\bm{SC}(t_f)$，$\bm{g}(t_f)=\bm{C}^T(t_f)\bm{Sz}(t_f)$
        \subsubsection{非零点调节器问题}
            对于线性定常系统，要求$\lim\limits_{t\rightarrow\infty}\bm{\tilde y}(t)=\lim\limits_{t\rightarrow\infty}(\bm{y}(t)-\bm{y}_0)=0$。
            稳态时，有$\bm{y}_0=\bm{Cx}_0$，$\bm{0}=\bm{Ax}_0+\bm{Bu_0}$，定义各偏差量$\bm{\tilde u}(t),\bm{\tilde x}(t),\bm{\tilde y}(t)$，它们同样满足状态方程和输出方程。
            
            由输出调节器求解得$\bm{\tilde u}(t)=-\bm{R^{-1}B^T P}(t)\bm{\tilde x}(t)=-\bm{K\tilde x}(t)$，即$\bm{u}(t)=-\bm{Kx}(t)+\bm{Kx}_0+\bm{u}_0=-\bm{Kx}(t)+\bm{g}$
        \subsubsection{PI跟踪控制器}
            对于线性定常系统，$\bm{\eta}$是要跟踪的常值向量，性能指标
            \[
            J=\dfrac12\int_{0}^{\infty}\left[(\bm{\eta-Cx})^T\bm{Q}(\bm{\eta-Cx})+\left(\dfrac{\d\bm{u}}{\d t}\right)^T\bm{R}\dfrac{\d\bm{u}}{\d t}\right]\d t
            \]

            引入$\bm{z}=\bm{\eta-Cx},\bm{V}(t)=\dfrac{\d\bm{u}}{\d t}$，引入增广状态变量$\bm{\bar x}=[\bm{\dot x}^T\quad\bm{z}^T]^T$，
            $\bm{\bar A}=\begin{bmatrix}\bm{A}& 0\\-\bm{C}& 0\end{bmatrix},\bm{\bar B}=\begin{bmatrix}\bm{B}\\0\end{bmatrix}$，
            
            则增广状态方程$\bm{\dot{\bar x}}=\bm{\bar A\bar x+\bar B V}$，
            得最优控制$\bm{V}^*(t)=-\bm{R}^{-1}[\bm{B}^T\quad 0]\bm{P\bar x}=-K_1\bm{\dot x}(t)-K_2\bm{z}(t)$，则$\bm{u}^*(t)=-K_1\bm{x}(t)+K_2\int(\bm{\eta-Cx}(t))\d t$
\newpage
\section{动态规划}
    动态规划是一种非线性规划，基于最优性原理，将复杂的优化问题变为易于求解的具有递推形式的多级决策问题。

    最优性原理：一个多级决策问题存在的最优策略具有这样的性质：不论初始级、初始状态和初始决策如何，当把其中任何一级和状态再用为初始级和初始状态时，余下的策略对此必定构成一个最优策略。

    不变嵌入原理：为解决一个特定的最优决策问题，而把原问题嵌入到一系列相似的易于求解的问题中去。对于多级决策过程来说，就是将多级决策问题化为一系列单级决策问题来处理。
    \subsection{离散型动态规划}
        离散系统的状态差分方程为$\bm{X}(k+1)=\bm{f}[\bm{X}(k),\bm{U}(k)]$，每步转移的指标泛函为$J=J[\bm{X}(k),\bm{U}(k)]$，其中$\bm{X}(k)$是$n\times1$的时刻k上的状态向量，$\bm{U}(k)$是$r\times1$的时刻k上的容许控制。

        N级决策过程，系统的状态向量序列是
        \begin{align*}
            \bm{X}(1)&=\bm{f}[\bm{X}(0),\bm{U}(0)]\\
            \bm{X}(2)&=\bm{f}[\bm{X}(1),\bm{U}(1)]=\bm{f}[\bm{f}[\bm{X}(0),\bm{U}(0)],\bm{U}(1)]\\
            \vdots&\\
            \bm{X}(N)&=\bm{f}[\bm{X}(N-1),\bm{U}(N-1)]=\bm{F}[\bm{X}(0),\bm{U}(0),\bm{U}(1),\cdots,\bm{U}(N-1)]
        \end{align*}
        相应的指标泛函为
        \begin{align*}
            J_1&=J[\bm{X}(0),\bm{U}(0)]\\
            J_2&=J[\bm{X}(0),\bm{U}(0)]+J[\bm{X}(1),\bm{U}(1)]\\
            \vdots&\\
            J_N&=\sum\limits_{k=0}^{N-1}J[\bm{X}(k),\bm{U}(k)]
        \end{align*}
        基于最优性原理，要求不论第一级控制向量$\bm{U}(0)$如何确定，其余N-1级过程必须构成N-1级最优控制过程。记$J_N^0[\bm{X}(0)]$为N级决策过程指标泛函的极小值，则可以得到递推函数方程如下：
        \[
        J_N^0[\bm{X}(0)]=\min\{J[\bm{X}(0),\bm{U}(0)]+J_{N-1}^0[\bm{f}[\bm{X}(0),\bm{U}(0)]]\}
        \]
        求解递推方程即可得到最优控制策略$\{\bm{U}^0(0),\bm{U}^0(1),\cdots,\bm{U}^0(N-1)\}$
    \subsection{连续型动态规划}
        指标泛函为$J=\Phi[\bm{X}(t_f),t_f]+\int_{t_0}^{t_f}L[\bm{X}(t),\bm{U}(t),t]\d t$

        假设末端时刻给定且最优控制向量确定，指标泛函的极小值仅是初始状态和初始时刻的标量函数
        \[
        J^0[\bm{X}(t),t]=J[\bm{X}(t_0),t_0]=\Phi[\bm{X}^0(t_f),t_f]+\int_{t_0}^{t_f}L[\bm{X}^0(t),\bm{U}^0(t),t]\d t
        \]

        基于最优性原理，$t$是区间$[t_0,t_f]$内一点，由$t$到$t_f$的过程分成$[t,t+\Delta t]$和$[t+\Delta t,t_f]$两级过程，这时指标泛函极小值
        \begin{align*}
        J^0[\bm{X}(t),t]&=\min\{\int_{t}^{t+\Delta t}L[\bm{X}(t),\bm{U}(t),t]\d t+\int_{t+\Delta t}^{t_f}L[\bm{X}(t),\bm{U}(t),t]\d t+\Phi[\bm{X}^0(t_f),t_f]\}\\
        &=\min\{\int_{t}^{t+\Delta t}L[\bm{X}(t),\bm{U}(t),t]\d t+J^0[\bm{X}(t+\Delta t),t+\Delta t]\}
        \end{align*}

        积分中值定理将第一项写为$\int_{t}^{t+\Delta t}L[\bm{X}(t),\bm{U}(t),t]\d t=L[\bm{X}(t),\bm{U}(t),t]\Delta t$，
        
        利用泰勒展开将第二项写为$J^0[\bm{X}(t+\Delta t),t+\Delta t]=J^0[\bm{X}(t),t]+\dfrac{\partial}{\partial\bm{X}^T}J^0[\bm{X}(t),t]\Delta\bm{X}+\dfrac{\partial}{\partial t}J^0[\bm{X}(t),t]\Delta t$

        带回上式可得贝尔曼（Bellman）方程：
        \begin{align*}
        \dfrac{\partial}{\partial t}J^0[\bm{X}(t),t]&=-\min\limits_{\bm{U}(t)}\{\bm{L}[\bm{X}(t),\bm{U}(t),t]+\dfrac{\partial}{\partial\bm{X}^T}J^0[\bm{X}(t),t]\bm{\dot X}(t)\}\\
        &=-\min\limits_{\bm{U}(t)}\{\bm{L}[\bm{X}(t),\bm{U}(t),t]+\dfrac{\partial}{\partial\bm{X}^T}J^0[\bm{X}(t),t]\bm{f}[\bm{X}(t),\bm{U}(t),t]\}
        \end{align*}

        右侧式子的$\min$部分对$\bm{U}^0(t)$求极小值并带回原式可得哈密顿-雅可比（Hamilton-Jacobi）方程：
        \[
        \dfrac{\partial J^0[\bm{X}^0(t),t]}{\partial t}+\bm{L}[\bm{X}^0(t),\bm{U}^0(t),t]+\dfrac{\partial J^0[\bm{X}^0(t),t]}{\partial\bm{X}^T(t)}\bm{f}[\bm{X}^0(t),\bm{U}^0(t),t]=0
        \]
\end{document} 